\setcounter{subsection}{2}

\subsection{Linear Maps}

\begin{appendixexercise}[3e11]
    Suppose $U = \set{(x_1, x_2, \ldots) \in \f^\infty \mid x_k \neq 0 \text{ for only finitely many } k}$.
    \begin{subexercises}
        \item Show that $U$ is a subspace of $\f^\infty$.
        \item Prove that $\f^\infty / U$ is infinite-dimensional.
    \end{subexercises}
\end{appendixexercise}

\begin{proof}
    \leavevmode
    \begin{subexercises}
        \item Since $(0,0,\dots) \in U$, $U$ is nonempty. Suppose $x$, $y \in U$ and $c \in \mathbb{F}$. Let $S_x$ and $S_y$ be the sets of indices where $x$ and $y$ have nonzero terms, respectively. By definition, $S_x$ and $S_y$ are finite. $cx+y$ can only have nonzero terms at indices in $S_x \cup S_y$. Since the union of two finite sets is finite, $cx+y \in U$, which makes $U$ a subspace of $\mathbb{F}^{\infty}$. 

        \item Partition $\mathbb{Z}^{+}$ into a countably infinite family of disjoint infinite subsets $A_1$, $A_2$, $\dots$. For each $n$, define $v_n \in \mathbb{F}^{\infty}$ such that its $j$th term is $1$ if $j \in A_n$ and $0$ otherwise. \\ 

        Suppose that $\sum^{m}_{i=1} c_i (v_i+U)=0+U$ for some positive integer $m$ and scalars $c_1$, $\dots$, $c_m \in \mathbb{F}$. Then, $w=\sum^{m}_{i=1} c_i v_i \in U$. By construction, for $i \in \{1, \dots, m\}$, the $j$-th term of $w$ is $c_i$ for all $j \in A_i$. Since $w \in U$, it only has a finite number of nonzero terms. But each $A_i$ is infinite, so this forces $c_i=0$ for all $i \in \{1, \dots, m\}$. Hence, $v_1+U$, $v_2+U$, $\dots$ are linearly independent in $\mathbb{F}^{\infty}/U$. \\ 

        Because $\mathbb{F}^{\infty}/U$ contains an infinite sequence of linearly independent vectors, it is infinite-dimensional.
    \end{subexercises}
\end{proof}

\begin{appendixexercise}[3d19]
    Suppose $V$ is finite-dimensional and $T \in \lin{V}$. Prove that $T$ has the same matrix with respect to every basis of $V$ if and only if $T$ is a scalar multiple of the identity operator.
\end{appendixexercise}

\begin{proof}[Coolempire93]
    The manual solution provides a constructive proof utilizing sums of the form $v_1 + v_k$ to eliminate the $k$-th column, along with $v_2 + v_1$ to eliminate the 1st column. Here, we present different constructions. Each proof assumes the declaration of the matrix to be $A$.
    \begin{enumerate}
        \item This method shows that we can swap two \emph{arbitrary} basis vectors to force the same value of all off-diagonal entries as well as the same value of all diagonal entries. Then, we can negate one basis vector to force all off-diagonal entries to be 0 via the linear independence of the basis vectors.
        
        Let $i,j$ be arbitrary with $1 \leq i,j \leq n$ and $i \neq j$. Consider the basis $w_1, \dots, w_n$ of $V$ defined by $w_i = v_j$, $w_j = v_i$ and $w_k = v_k$ otherwise. Then 
        \begin{align*}
        T(w_j - v_i) &= Tw_j - Tv_i \\
        T(0) &= A_{1,j}w_1 + A_{2,j}w_2 + \cdots + A_{n,j}w_n - (A_{1,i}v_1 + A_{2,i}v_2 + \cdots + A_{n,i}v_n) \\
        0 &= (A_{j,j} - A_{i,i})v_i + (A_{i,j} - A_{j,i})v_j + \sum_{\substack{k \neq i \\ k \neq j}}(A_{k,j} - A_{k,i})v_k.
        \end{align*}
        By linear independence of $v_1, \dots, v_n$, we have $A_{i,i} = A_{j,j}$, $A_{i,j} = A_{j,i}$ and $A_{k,j} = A_{k,i}$ for all $k \neq i,j$.
        Because $i$ and $j$ were arbitrary, we have $A_{1,1} = A_{2,2} = \cdots = A_{n,n} = \lambda$, as well as $A$ is symmetric, and all off-diagonal entries match across rows, so we can say $A_{i,j} = \gamma$ for all $i \neq j$.
        To finish, consider the basis $u_1, \dots, u_n$ defined by $u_1 = -v_1$ and $u_k = v_k$ otherwise. Then 
        \begin{align*}
        T(v_1 + u_1) &= Tv_1 + Tu_1 \\
        T(0) &= \lambda v_1 + \gamma(v_2 + \cdots + v_n) + (\lambda u_1 + \gamma(u_2 + \cdots + u_n)) \\
        0 &= 2\gamma(v_2 + \cdots + v_n).
        \end{align*}
        Again by linear independence of $v_1, \dots, v_n$, we have $2\gamma = 0$ and thus $\gamma = 0$. Therefore, $A = \lambda I$ and $T$ is a scalar multiple of the identity.
        \item This method shows that we can negate one \emph{arbitrary} basis vector to force all off-diagonal entries to be 0 via the linear independence of the basis vectors. Then, we can swap $v_1$ with any other basis vector to force all diagonal entries to be the same.
        
        Let $j$ be arbitrary and consider the basis $w_1, \dots, w_n$ of $V$ defined by $w_j = -v_j$ and $w_i = v_i$ otherwise. Then
        \begin{align*}
        T(w_j + v_j) &= Tw_j + Tv_j \\
        T(0) &= A_{1,j}w_1 + A_{2,j}w_2 + \cdots + A_{n,j}w_n + A_{1,j}v_1 + A_{2,j}v_2 + \cdots + A_{n,j}v_n \\
        0 &= 2\sum_{k \neq j}A_{k,j}v_k.
        \end{align*}
        By linear independence of $v_1, \dots, v_n$, we have $A_{k,j} = 0$ for all $k \neq j$. Since $j$ was arbitrary, we have that all off-diagonal entries of $A$ are 0.

        We turn to one more basis to prove that all diagonal entries of $A$ must be the same. Once again, let $j$ be arbitrary, and this time consider the basis $u_1, \dots, u_n$ defined by $u_1 = v_j$, $u_j = v_1$ and $u_i = v_i$ otherwise. Then
        \begin{align*}
        T(v_1 - u_j) &= Tv_1 - u_j \\
        T(0) &= A_{1,1}v_1 - A_{j,j}u_j \\
        0 &= (A_{1,1} - A_{j,j})v_1.
        \end{align*}
        As $v_1$ belongs to a basis for $V$, it is nonzero, and thus as $j$ was arbitrary, $A_{1,1} = A_{2,2} = \cdots = A_{n,n} = \lambda$ for some scalar $\lambda$, and we have $A = \lambda I$, i.e., $T$ is a multiple of the identity operator. \qh
    \end{enumerate}
\end{proof}